\section*{Erd\H{o}s Problem 1130: the maximin interpolation problem}

\paragraph{Statement and normalization.}
For distinct ordered nodes \(-1\leq x_1<\cdots<x_n\leq1\), let
\[
 L_{\boldsymbol x}(x)=\sum_{k=1}^n
 \left|\prod_{j\ne k}\frac{x-x_j}{x_k-x_j}\right|,
 \qquad
 \Upsilon(\boldsymbol x)=\min_{0\leq i\leq n}
 \max_{x\in[x_i,x_{i+1}]}L_{\boldsymbol x}(x),
\]
where \(x_0=-1,x_{n+1}=1\). We prove the logarithmic bound and characterize all maximizers under this literal convention, which includes both exterior intervals.
Source: \url{https://www.erdosproblems.com/1130}.

\paragraph{The canonical input.}
Theorems 1 and 2 of de Boor--Pinkus (1978), \url{https://doi.org/10.1016/0021-9045(78)90014-X}, imply that for \(n\geq3\) there is a unique canonical array
\(-1=t_1^*<\cdots<t_n^*=1\) whose \(n-1\) internal interval maxima all have a common value \(\Lambda_n^*\).
For any other canonical array, the minimum internal maximum is strictly smaller and the maximum internal maximum is strictly larger.
We use this theorem as an explicit external input, with its endpoints fixed and its node indexing adjusted. Write \(L^*=L_{\boldsymbol t^*}\).

\paragraph{Characterization with exterior intervals included.}
The exact maximum of \(\Upsilon\) is \(\Lambda_n^*\), and all maximizers are
\[
 \boxed{\quad x_i=c+h t_i^*,\quad h>0,\quad
 -1\leq c-h<c+h\leq1,\quad
 L^*((-1-c)/h)\geq\Lambda_n^*,\
 L^*((1-c)/h)\geq\Lambda_n^*.\quad}
\]
Indeed, affine normalization by \(c=(x_1+x_n)/2\), \(h=(x_n-x_1)/2\) preserves every internal maximum. Thus
\(\Upsilon(\boldsymbol x)\leq\min_{1\leq i<n}\lambda_i\leq\Lambda_n^*\),
with equality only if the normalized array is \(\boldsymbol t^*\).
For this array each exterior maximum occurs at its outer endpoint: beyond the outermost node each absolute Lagrange product grows as one moves farther away. The two displayed inequalities are therefore exactly the remaining equality conditions.

The canonical array is symmetric by uniqueness. Its exterior function grows strictly from \(L^*(1)=1\) to infinity, and \(\Lambda_n^*>1\).
Let \(r_n>1\) satisfy \(L^*(r_n)=\Lambda_n^*\). The conditions become
\[
 |c|<1,\qquad 0<h\leq(1-|c|)/r_n.
\]
They are nonempty, proving attainment, and also show that maximizers need not be unique or have all \(n+1\) maxima equal. For example, with three nodes \(-t,0,t\), the internal maxima are \(5/4\), the exterior ones are \(2/t^2-1\), and every \(0<t\leq\sqrt{8/9}\) maximizes \(\Upsilon\).
For \(n=2\), \(\Upsilon=1\) for every pair of nodes, because its internal maximum is \(1\) and all maxima are at least \(1\). For \(n=1\), likewise \(\Upsilon=1\).

\paragraph{An elementary \(O(\log n)\) upper bound.}
It remains to bound \(\Lambda_n^*\). Use the \(n\) Chebyshev zeros
\(z_k=\cos\theta_k\), \(\theta_k=(2k-1)\pi/(2n)\).
For \(x=\cos\theta\), their Lagrange functions satisfy
\[
 |l_k(\cos\theta)|
 =\frac{|\cos n\theta|\sin\theta_k}
 {n|\cos\theta-\cos\theta_k|}
 \leq\frac{\pi|\cos n\theta|}{n|\theta-\theta_k|}.
\]
The inequality follows from
\(\sin\theta_k\leq2\sin((\theta+\theta_k)/2)\) and
\(\sin(|\theta-\theta_k|/2)\geq|\theta-\theta_k|/\pi\).
Let \(k_0\) be a nearest zero angle. Then
\(|\cos n\theta|\leq n|\theta-\theta_{k_0}|\), so the \(k_0\)-term is at most \(\pi\). For \(k\ne k_0\),
\[
 |\theta-\theta_k|\geq(|k-k_0|-1/2)\pi/n,
 \qquad |l_k(\cos\theta)|\leq\frac1{|k-k_0|-1/2}.
\]
Consequently the sum is at most
\(\pi+2\sum_{j=1}^{n-1}(j-1/2)^{-1}=O(\log n)\), uniformly in \(\theta\); values at nodes follow by continuity.
Stretching these zeros so that the extreme nodes become \(-1,1\) cannot increase the maximum: the affine preimage of the domain is a subinterval of the original domain.
Canonical optimality now gives \(\Lambda_n^*=O(\log n)\), and hence \(\Upsilon\ll\log n\) for \(n\geq2\).

\paragraph{Completion estimate.}
The logarithmic bound and complete literal maximizer characterization follow, with only the canonical extremal-node theorem imported.
\textbf{COMPLETION: 100\%}
